\documentclass[11pt]{article}
\usepackage{coursenotes}
\begin{document}
\handoutheader{H1}{Interference and marketplace experiments}{When one unit's treatment changes another's outcome}{Meeting 2 (Experiments)}

\begin{decision}
A platform tests a delivery-fee discount on a random half of customers and sees treated customers order 33\%
more. If couriers are scarce, part of that gain is taken from control customers, whose orders now wait longer.
The number the rollout decision needs, ``everyone discounted versus no one'', may be far smaller, and even zero.
The questions are what an experiment estimates when units interfere, and which designs recover the effect a
rollout needs.
\end{decision}

\begin{goals}
  \item write potential outcomes that depend on others' treatments, using an exposure mapping;
  \item distinguish direct, spillover, and total (global) effects, and show what a unit-randomized test estimates;
  \item explain marketplace bias from shared supply, and compute it in a simple model;
  \item choose among cluster, two-stage, switchback, and two-sided designs, and state what each assumes;
  \item analyze each design with the right unit of inference.
\end{goals}

%=====================================================================================
\sectionone

\subsection{Potential outcomes under interference}

Without SUTVA (the assumption that each unit's outcome depends only on its own treatment), unit $i$'s outcome
depends on the whole assignment vector $\mathbf{d} = (d_1, \dots, d_n)$, written $Y_i(\mathbf{d})$. With $2^n$ potential outcomes per unit nothing is estimable without structure. An
\emph{exposure mapping} supplies it: assume $Y_i(\mathbf{d}) = Y_i(d_i, h_i)$, where $h_i = h_i(\mathbf{d})$ summarizes the
treatment of the units that can affect $i$, such as the share treated among $i$'s friends, in $i$'s cluster, or in
$i$'s city during the same hour.

\begin{defn}{Effects under interference}{effects}
For saturations $h, h' \in [0,1]$:
\begin{itemize}[leftmargin=1.2em]
  \item direct effect at saturation $h$: $\text{DE}(h) = \E[Y(1,h) - Y(0,h)]$;
  \item spillover effect on arm $d$: $\text{SE}_d(h, h') = \E[Y(d,h) - Y(d,h')]$;
  \item total (global) treatment effect: $\text{GTE} = \E[Y(1,1) - Y(0,0)]$, everyone treated versus no one.
\end{itemize}
\end{defn}

A rollout decision needs the GTE (or its net-value version). A unit-randomized test delivers something else.

\begin{prop}{What a unit-randomized test estimates}{naive}
Suppose a share $p$ of units is randomized to treatment and each unit's exposure is $h_i = p$ (every unit
experiences the same treated share, as in a well-mixed market). Then the difference in means estimates
$\text{DE}(p)$, and
\[
  \text{GTE} - \text{DE}(p) = \underbrace{\E[Y(1,1) - Y(1,p)]}_{\text{treated units' spillover}} - \underbrace{\E[Y(0,0) - Y(0,p)]}_{\text{control units' spillover}} .
\]
\end{prop}
\begin{proof}
Randomization makes $\E[Y \mid D=d] = \E[Y(d,p)]$, so the difference is $\text{DE}(p)$. Add and subtract
$\E[Y(1,p)]$ and $\E[Y(0,p)]$ in $\text{GTE} = \E[Y(1,1)] - \E[Y(0,0)]$.
\end{proof}

The sign of the gap depends on the mechanism. \emph{Competition} for a shared resource (couriers, inventory, ad
slots, a seller's attention) makes control units worse off when others are treated, so $\E[Y(0,p)] < \E[Y(0,0)]$ and
the test \emph{overstates} the GTE. \emph{Network effects} (a referral program, a social feature) make controls
better off, and the test \emph{understates} it.

\subsection{Marketplace bias in a simple model}

\begin{prop}{Shared supply}{supply}
Let each customer request $q_1$ orders if treated and $q_0 < q_1$ if not, and let couriers complete at most $K$ orders
per customer on average, allocated in proportion to requests. If $q_0 \le K < q_1$ and a share $p$ is treated, the
completion rate is $c(p) = \min\{1, K/(pq_1 + (1-p)q_0)\}$, the unit-randomized test estimates
$\text{DE}(p) = (q_1 - q_0)\,c(p)$, and the GTE is $K - q_0$.
\end{prop}
\begin{proof}
Each customer's completed orders are her requests times the common completion rate, so arm means are $q_1c(p)$ and
$q_0c(p)$. With everyone treated, requests $q_1 > K$ are capped at $K$; with no one treated, $q_0 \le K$ are all
completed.
\end{proof}

When $K = q_0$, supply is exactly used up already and the GTE is zero: the discount only moves orders from some
customers to others. The test still shows a positive effect at every $p$, because treated customers win the
competition for couriers. Larger tests do not help; only the design does.

\subsection{Designs that recover the total effect}

\paragraph{Cluster randomization.} Partition units into clusters (cities, delivery zones, graph communities) such
that interference stays within a cluster, and randomize whole clusters.

\begin{prop}{Cluster randomization identifies the GTE}{cluster}
If each unit's outcome depends only on the assignments within its own cluster, and clusters are assigned entirely to
treatment or control at random, the difference in cluster-level mean outcomes is unbiased for the GTE (averaged over
clusters).
\end{prop}
\begin{proof}
In a treated cluster every unit has $(d_i, h_i) = (1,1)$; in a control cluster $(0,0)$. Clusters are the units of a
completely randomized experiment with outcomes $\bar Y_g(1,1)$ and $\bar Y_g(0,0)$, and under complete randomization
the difference in means is unbiased for the average effect over the units randomized.
\end{proof}

The trade-off is bias against variance: bigger clusters contain more of the interference but leave fewer
independent units, and the design effect (the inflation of variance caused by correlated outcomes within a
cluster, which grows with cluster size) widens the interval. Interference that crosses cluster
borders (customers near a zone boundary) biases the estimate toward the unit-level answer.

\paragraph{Two-stage (randomized saturation) designs.} First randomize each cluster's saturation $h$ (say 0\%, 20\%,
50\%, 80\%), then randomize units within clusters at that rate. Comparing treated and control units within clusters
estimates $\text{DE}(h)$ at each $h$; comparing controls across saturations estimates $\text{SE}_0(h, h')$. The full curve
lets the platform predict a rollout at any share, not only at 100\%.

\paragraph{Switchback designs.} When the whole market interacts (one city's couriers serve everyone), randomize
\emph{time}: each block (an hour, a day) is treated or not, for the whole market. Treatment in one block can carry over
into the next (orders queued, couriers repositioned). If carryover lasts at most $m$ periods, dropping the first $m$
periods of each block (a \emph{washout}) leaves blocks whose outcomes depend only on their own assignment, and the
block-level difference in means is unbiased. Longer blocks reduce carryover bias but leave fewer blocks; time-of-day
patterns call for randomizing within days or balancing blocks across hours.

\paragraph{Two-sided randomization.} In a market with buyers and sellers, randomize both sides and treat a transaction
only when both parties are treated (or apply the treatment by the transaction's side). Comparing the four cells
separates demand-side from supply-side competition; Johari et al.\ (2022) show it can reduce bias from both at
moderate cost in variance.

\paragraph{Budget-split designs.} For advertising with budgets, split each advertiser's budget into two independent
halves, one per arm, so the arms do not compete for the same budget. This removes cannibalization through the
budget, the dominant channel in ad marketplaces.

\subsection{Inference}

The unit of inference is the unit of randomization: clusters for cluster designs (cluster-robust SEs or cluster-level
means), blocks for switchbacks (randomization inference over block assignments is natural, since there are few blocks
and outcomes are serially correlated). Two-stage designs need SEs at the cluster level for spillover contrasts.
Interference also invalidates the usual SE of a unit-randomized test, because units' outcomes are correlated through
the shared resource.

\begin{pitfall}[reading a marketplace test]
Treating a unit-randomized lift as the rollout effect when units share supply; defining clusters by convenience
rather than by where interference stops; analyzing a switchback with hour-level SEs as if hours were independent.
\end{pitfall}

%=====================================================================================
\sectiontwo

\paragraph{Setting.} Pujiang Delivery tests a delivery-fee discount in one city. Customers request $0.30$ orders a week on average
without the discount and $0.40$ with it. Courier capacity completes at most $0.30$ orders per customer per week, and
dispatch serves requests in proportion to their number. Each completed order earns ¥8 of margin; the discount costs
¥3 per completed discounted order.

\paragraph{Step 1: the user-randomized test (50/50).} Requests average $0.35$, so the completion rate is
$0.30/0.35 = 0.857$. Treated customers complete $0.40 \times 0.857 = 0.343$ orders a week and controls
$0.30 \times 0.857 = 0.257$, an estimated lift of $0.086$, or 33\%. The net per treated customer is
$8 \times 0.086 - 3 \times 0.343 = -$¥0.34. The test loses money even taken at face value, and the face value is itself
misleading.

\paragraph{Step 2: the total effect.} With everyone discounted, requests of $0.40$ are capped at $0.30$; with no one
discounted, $0.30$ are all completed. $\text{GTE} = 0$ (Result~\ref{prop:supply}). A citywide discount buys no extra
orders and costs $3 \times 0.30 = $ ¥0.90 per customer per week.

\paragraph{Step 3: a two-stage design.} Running saturations of 20\% and 80\% in other cities, alongside the 50\% test:
\begin{center}
\begin{tabular}{lcccc}
\toprule
Saturation $h$ & Completion rate & Treated orders & Control orders & $\text{DE}(h)$ \\
\midrule
20\% & 0.938 & 0.375 & 0.281 & 0.094 \\
50\% & 0.857 & 0.343 & 0.257 & 0.086 \\
80\% & 0.789 & 0.316 & 0.237 & 0.079 \\
\bottomrule
\end{tabular}
\end{center}
The direct effect barely moves, but controls lose orders as saturation rises ($0.281 \to 0.237$). That is the
spillover the unit-level test cannot see. Extrapolating the control curve to $h = 0$ gives $0.30$, and the treated curve to $h = 1$
gives $0.30$, so the GTE is zero.

\paragraph{Step 4: a switchback.} Pujiang Delivery instead alternates the discount citywide by day for six weeks, with the first
two hours of each day dropped as washout (orders queued overnight). Treated and control days complete the same number
of orders, confirming that supply binds. The recommendation is to spend the discount budget on courier supply in
peak hours, where it can raise completed orders, and to re-test there.

\begin{exercise}[computation]
Suppose courier capacity rises to $0.35$ orders per customer. Recompute the 50/50 test's estimated lift and the GTE.
\end{exercise}
\answer{Requests average $0.35 = K$, so completion is $1$: treated $0.40$, control $0.30$, estimated lift $0.10$. With
everyone treated, requests $0.40$ are capped at $0.35$, so $\text{GTE} = 0.35 - 0.30 = 0.05$, half the test's estimate.}

\begin{exercise}[judgment]
The growth team proposes randomizing by customer but ``only analyzing customers in zones with spare courier
capacity''. When does this recover the GTE, and what would you need to check?
\end{exercise}
\answer{It recovers the GTE in those zones only if capacity stays slack under a full rollout (so no competition) and
the spare-capacity zones are chosen before outcomes and without reference to treatment. Check utilization under the
test, whether treatment itself changes which zones look slack (a post-treatment selection), and whether those zones
resemble the rollout population. The result transports only to zones that would also remain slack.}

%=====================================================================================
\sectionthree

\begin{extension}[Equilibrium effects and pricing]
When treatment changes prices that clear the market (surge pricing, courier pay), effects pass through equilibrium.
Wager and Xu (2021) show how small randomized perturbations to each unit's price can estimate marginal equilibrium
effects without large-scale tests.
\end{extension}

\subsection*{Annotated reading}
\begin{readinglist}
\reading{Hudgens, Michael G., and M. Elizabeth Halloran. 2008. ``Toward Causal Inference with Interference.''
\emph{Journal of the American Statistical Association} 103 (482): 832--842.
\url{https://doi.org/10.1198/016214508000000292}}{Direct, indirect, and total effects with two-stage
designs}{Sections 1--4}{3}
\reading{Aronow, Peter M., and Cyrus Samii. 2017. ``Estimating Average Causal Effects under General
Interference, with Application to a Social Network Experiment.'' \emph{Annals of Applied Statistics} 11 (4):
1912--1947. \url{https://doi.org/10.1214/16-AOAS1005}}{Exposure mappings and design-based
estimators}{Sections 1--3}{3}
\reading{Johari, Ramesh, Hannah Li, Inessa Liskovich, and Gabriel Y. Weintraub. 2022. ``Experimental Design
in Two-Sided Platforms: An Analysis of Bias.'' \emph{Management Science} 68 (10): 7069--7089.
\url{https://doi.org/10.1287/mnsc.2021.4247}}{A marketplace model of when customer- or listing-side tests are
biased, and two-sided randomization}{Sections 1--5}{2}

\reading{Bojinov, Iavor, David Simchi-Levi, and Jinglong Zhao. 2023. ``Design and Analysis of Switchback
Experiments.'' \emph{Management Science} 69 (7): 3759--3777.
\url{https://doi.org/10.1287/mnsc.2022.4583}}{Optimal block lengths under carryover}{Sections 1--4}{3}
\reading{Holtz, David, Felipe Lobel, Ruben Lobel, Inessa Liskovich, and Sinan Aral. 2025. ``Reducing
Interference Bias in Online Marketplace Experiments Using Cluster Randomization: Evidence from a Pricing
Meta-experiment on Airbnb.'' \emph{Management Science} 71 (1): 390--406.
\url{https://doi.org/10.1287/mnsc.2020.01157}}{Evidence from a pricing meta-experiment: how much bias
clusters remove}{All}{1}
\reading{Ugander, Johan, Brian Karrer, Lars Backstrom, and Jon Kleinberg. 2013. ``Graph Cluster Randomization:
Network Exposure to Multiple Universes.'' In \emph{Proceedings of the 19th ACM SIGKDD International Conference
on Knowledge Discovery and Data Mining}, 329--337. arXiv:1305.6979.}{Designing clusters on a social
network}{Sections 1--3}{2}
\reading{Wager, Stefan, and Kuang Xu. 2021. ``Experimenting in Equilibrium.'' \emph{Management Science} 67
(11): 6694--6715. \url{https://doi.org/10.1287/mnsc.2020.3844}}{Local experimentation for equilibrium effects
in marketplaces}{Sections 1--3}{3}
\end{readinglist}

\end{document}
